\documentclass[../../../include/open-logic-section]{subfiles}

\begin{document}

\olfileid[fr]{sth}{replacement}{absinf}
\olsection{Le remplacement et l'«~infini absolu~»}

Laissons maintenant \limofsize{} de côté pour explorer une autre
famille de tentatives de justification (intrinsèque) du remplacement.
Celles-ci prennent au sérieux l'idée que les ensembles se forment
par étapes.

Lorsque nous avons esquissé pour la première fois le processus
itératif, nous avons proposé des principes qui expliquaient ce
qui se passe à chaque étape : \stageshier, \stagesord{} et
\stagesacc. Plus tard, nous avons ajouté des principes portant
sur le nombre d'étapes : \stagessucc{} disait que le processus
de formation des ensembles ne s'arrête jamais, et \stagesinf{}
qu'il atteint une étape d'indice infini.

On peut raisonnablement suggérer que ces deux derniers principes
découlent d'un principe plus général, tel que :
\begin{enumerate}
	\item[] \stagesinex. Il y a une infinité absolue d'étapes ;
	la hiérarchie est aussi haute qu'elle pourrait l'être.
\end{enumerate}
Il s'agit évidemment d'un principe informel. Mais, même s'il ne
\emph{découle} pas immédiatement de la conception cumulative et
itérative des ensembles, il paraît certainement \emph{compatible}
avec elle. À tout le moins, et à la différence de \limofsize,
il conserve l'idée que les ensembles se forment étape par étape.

Nous espérons maintenant tirer de \stagesinex{} une justification
du remplacement. Voyons comment.

Dans la \olref[ordinals][idea]{sec}, nous avons vu qu'il est facile
de construire un bon ordre qui devrait, intuitivement, être
isomorphe à $\omega+\omega$. Autrement dit, nous pouvons aisément
imaginer un processus itératif étape par étape dont le type
d'ordre serait, intuitivement, $\omega+\omega$. Si nous avons
accepté \stagesinex, nous devrions donc sûrement admettre
qu'il y a au moins une étape $\omega+\omega$-ième de la
hiérarchie, c'est-à-dire $V_{\omega+\omega}$ : la hiérarchie
\emph{pourrait} certainement aller jusque-là.

Cette idée se généralise : pour tout bon ordre, la construction
de la hiérarchie itérative devrait se poursuivre au moins aussi
loin que ce bon ordre. Il suffirait, pour le garantir, de traiter
le \olref[ordinals][ordtype]{thmOrdinalRepresentation} comme un
\emph{axiome}. Celui-ci nous dirait que tout bon ordre est
isomorphe à un ordinal de von Neumann. Puisque chaque ordinal
de von Neumann est égal à son propre rang,
le \olref[ordinals][ordtype]{thmOrdinalRepresentation} nous
dirait alors que, chaque fois que nous pouvons décrire un bon
ordre dans notre théorie des ensembles, le processus itératif
de formation des ensembles doit dépasser ce bon ordre.

Cette idée paraît certainement découler de \stagesinex.
Malheureusement, si notre but est d'en déduire le remplacement,
un obstacle technique simple se présente : le remplacement est
strictement plus fort que le
\olref[ordinals][ordtype]{thmOrdinalRepresentation}.
(C'est ce que remarque \citet[\S13.2]{Potter2004} ; nous le
démontrerons dans la
\olref[sth][replacement][finiteaxiomatizability]{sec}.)

Ainsi, si nous voulons comprendre \stagesinex{} d'une manière
qui permette d'obtenir le remplacement, ce principe ne peut
\emph{pas se contenter} de dire que la hiérarchie dépasse tout
bon ordre. Il doit affirmer davantage. À cette fin,
\cite{Shoenfield:AST} a proposé un renforcement très naturel :
aucun ensemble n'est \emph{cofinal} dans la hiérarchie.\footnote{
G\"odel semble avoir proposé une idée voisine ; voir
\citet[p.~223]{Potter2004}. Pour une discussion de G\"odel
et de \citeauthor{Shoenfield:AST}, voir
\citet[90--5]{Incurvati2020}.}
Plus précisément, si $\tau$ est une opération définissable qui, à des
ensembles, associe des étapes de la hiérarchie, l'image de tout
ensemble $A$ par $\tau$ n'épuise pas la hiérarchie. Autrement
dit (de manière schématique) :
\begin{enumerate}
	\item[] \stagescofin. Si $A$ est un ensemble et si
	$\tau(x)$ est une étape pour tout $x \in A$, alors il
	existe une étape qui vient après chacune des $\tau(x)$,
	pour $x \in A$.
\end{enumerate}
Il est évident que $\ZF$ démontre une version convenablement
formalisée de \stagescofin. Inversement, on peut soutenir
informellement que \stagescofin{} justifie le remplacement.
\footnote{Il serait plus difficile de démontrer le remplacement
à partir d'une formalisation de \stagescofin, car $\Z$ seule
n'est pas assez forte pour définir les étapes ; la manière
de formaliser \stagescofin{} n'est donc pas claire. Une possibilité
serait toutefois de travailler dans une extension de $\LT$,
comme nous l'avons évoqué dans la
\olref[sth][replacement][extrinsic]{sec}.}
Supposons en effet que
$(\forall x \in A)\lexists![y][\phi(x,y)]$. Pour chaque
$x \in A$, soit $\sigma(x)$ l'unique $y$ tel que $\phi(x,y)$,
et soit $\tau(x)$ l'étape où $\sigma(x)$ se forme pour la
première fois. D'après \stagescofin, il existe une étape $V$
telle que $(\forall x \in A)\tau(x)\in V$. Puisque chaque
$\tau(x) \in V$ et que $\sigma(x) \subseteq \tau(x)$,
la séparation permet d'obtenir
$\Setabs{y \in V}{(\exists x \in A)\sigma(x) = y}
= \Setabs{y}{(\exists x \in A)\phi(x,y)}$.

\begin{prob}
	Formalisez \stagescofin{} dans $\ZF$.
\end{prob}

Ainsi, \stagescofin{} justifie le remplacement. Il est au moins
plausible que \stagesinex{} justifie à son tour \stagescofin.
Supposons en effet que \stagescofin{} soit faux. La hiérarchie
admet alors un ensemble cofinal~$A$ : il existe une opération définissable
$\tau$ telle que, pour toute étape~$S$, il y ait un $x \in A$
tel que $S \in \tau(x)$. Dans ce cas, nous disposons bien
d'un moyen de saisir le prétendu «~infini absolu~» de la
hiérarchie : il est \emph{épuisé} par l'image de $A$ sous
$\tau$. Cela compromet l'idée que la hiérarchie est
«~absolument infinie~». Par contraposition, \stagesinex{}
entraîne \stagescofin, qui justifie à son tour le remplacement.

Voilà une tentative réellement prometteuse de fournir une
justification intrinsèque du remplacement. Mais c'est à vous
de décider si elle réussit en définitive.

\end{document}
